\documentclass[11pt]{article}
\usepackage{amsmath,amssymb,amsthm,hyperref}
\newtheorem{theorem}{Theorem}
\title{Factorially Many Reciprocal Sheets in a Compact Positive Chart}
\author{}
\date{}
\begin{document}
\maketitle

\paragraph{Status and purpose.}
This is a complete written construction, not a Lean declaration.
It shows that far-away poles do not make reciprocal score hypersurfaces
a pseudohyperplane arrangement. It does not prove superlinear
shattering, and it supplies no universal head bound.

\begin{theorem}[Regular permutation sheets]
For every integer $H\ge1$ and every $0<\varepsilon<1/4$, there exist
$H$ rational functions of $b=(b_1,\ldots,b_H)\in(0,\varepsilon)^H$,
\[
 F_i(b)=\frac1{s_i}\sum_{h=1}^H\frac1{d_i-b_h}-1,
 \qquad 1\le i\le H,
\]
with $1-\varepsilon<d_i<1$ and $s_i>0$, whose common zero set
contains at least $H!$ distinct regular isolated points. All denominators
lie in $(1-2\varepsilon,1)$ throughout the parameter box.
\end{theorem}

\begin{proof}
Choose distinct numbers $\lambda_1,\ldots,\lambda_H$ in
$(0,\varepsilon)$ and distinct numbers $d_1,\ldots,d_H$ in
$(1-\varepsilon,1)$. Set
\[
 s_i=\sum_{j=1}^H\frac1{d_i-\lambda_j}>0.
\]
For every permutation $\pi$ of $\{1,\ldots,H\}$, the point
$b_h=\lambda_{\pi(h)}$ is a common zero, because summing over a
permutation leaves $s_i$ unchanged. These $H!$ points are distinct.
The Jacobian at such a point is
\[
 \frac{\partial F_i}{\partial b_h}
 =\frac1{s_i(d_i-\lambda_{\pi(h)})^2}.
\]
Let $C_{ij}=1/(d_i-\lambda_j)$. Cauchy's determinant formula makes
$\det C\ne0$, since the two sets of arguments are separately distinct
and all pairwise denominators are nonzero. Every entry of $C$ is
positive, so its permanent is strictly positive. Borchardt's identity
is the exact formula
\[
 \det(C\mathbin{\circ}C)=\det(C)\operatorname{perm}(C).
\]
Here $\circ$ denotes entrywise multiplication. A primary proof is
D.~Singer, \emph{A Bijective Proof of Borchardt's Identity},
Electronic Journal of Combinatorics 11 (2004), Theorem~1.1,
\url{https://www.combinatorics.org/ojs/index.php/eljc/article/download/v11i1r48/pdf/}.
Thus the squared Cauchy matrix is nonsingular. Row scaling by
$1/s_i$ and column permutation preserve nonsingularity, proving that
every displayed common zero is regular. The inverse function theorem
then makes each one isolated. Finally,
$1-2\varepsilon<d_i-b_h<1$ throughout the box.
\end{proof}

\paragraph{Compatibility with whole-cube admissibility.}
Fix any desired number $S$ of one-hot column variables, and choose
$\varepsilon<1/[4(H+S)]$. For a Boolean input
$(u,v)\in\{0,1\}^H\times\{0,1\}^S$, use the $H$ denominators
\[
 B_h(u,v)=1-\sum_{i=1}^H(1-d_i)u_i
             -\sum_{j=1}^S b_{hj}v_j,
 \qquad 0<b_{hj}<\varepsilon.
\]
Every coordinate slope is strictly negative, and on the entire cube
\[
 B_h(u,v)>1-(H+S)\varepsilon>\frac34.
\]
These are admissible negative-orientation affine denominators, uniformly
for every choice of all column parameters in the indicated box.
Choose the affine numerator
$A_h(u,v)=\sum_i u_i/s_i$, the same in every head, and output bias
$-1$. On the promised input $(u,v)=(e_i,e_j)$ the score is exactly
\[
 -1+\sum_{h=1}^H\frac{1/s_i}{d_i-b_{hj}}=F_i(b_{1j},\ldots,b_{Hj}).
\]
Thus the regular-sheet construction occurs inside the actual admissible
affine-ratio chart, with poles bounded away from all Boolean inputs.
Only the description of the selected rows uses the one-hot promise;
positivity and denominator orientation hold on every input.

\paragraph{Why this is not a shattering construction.}
The $H$ functions $F_i$ are symmetric in the head parameters $b_h$.
Every other query with a common row baseline and equal head coefficients
is symmetric as well, so it has the same value at every permutation
sheet. Factorially many parameter solutions therefore do not themselves
provide even one extra readable bit.

Asymmetric query rows can distinguish the sheets. For example,
\[
 \frac1{d-b_h}-\frac1{d-b_k}
 =\frac{b_h-b_k}{(d-b_h)(d-b_k)}
\]
reads a pairwise comparison whenever $d>\varepsilon$. But pairwise
comparisons of a permutation shatter at most $H-1$ fixed queries:
if the queried undirected graph has a cycle, orient that cycle
cyclically to obtain an impossible strict ordering; if it is a forest,
every edge orientation extends to a total order, so every labeling is
possible. A forest has at most $H-1$ edges. This is an elementary
obstruction to using pairwise comparisons as the missing
$\Theta(H\log H)$ decoder.

More generally, suppose a query differs from a symmetric reciprocal sum
at only one slot $g$. On the permutation sheets it has the form
\[
 Q(\pi)=c+\sum_{h\ne g}\frac{a}{d-\lambda_{\pi(h)}}
                 +\frac{a_g}{e-\lambda_{\pi(g)}}
 =c+a\sum_{j=1}^H\frac1{d-\lambda_j}
       -\frac{a}{d-t}+\frac{a_g}{e-t},
 \quad t=\lambda_{\pi(g)},
\]
where $d,e>\varepsilon$. Multiplication by the strictly positive
$(d-t)(e-t)$ gives a polynomial of degree at most two. Unless it is
identically zero, the query has at most two zero locations on
$(0,\varepsilon)$. Consequently, $r$ such queries attached to the same
slot have at most $2r+1$ strict sign vectors as $t$ varies. If they are
shattered, $2^r\le2r+1$, which forces $r\le2$. Partitioning a set of
one-slot queries by their slot proves that this restricted query class
cannot shatter more than $2H$ rows. This argument allows the symmetric
baseline, both coefficients, both poles, and the bias to vary from
query to query. It is a statement about the exact permutation sheets;
it does not control arbitrary neighborhoods of them.

There is also a concrete obstruction for queries which alter every
slot. Suppose $H\ge3$, $d_h>\varepsilon$, and
\[
 Q(b)=c+\sum_{h=1}^H\frac{\alpha_h+\beta_h b_h}{d_h-b_h}.
\]
Then $Q(\lambda_{\pi(1)},\ldots,\lambda_{\pi(H)})$ cannot have
strict sign $\operatorname{sgn}(\pi)$ for every permutation $\pi$.
Indeed, clearing the positive product of denominators gives the
multi-affine polynomial
\[
 P(b)=c\prod_h(d_h-b_h)
       +\sum_h(\alpha_h+\beta_h b_h)\prod_{k\ne h}(d_k-b_k).
\]
Every squarefree monomial $\prod_{h\in T}b_h$ has zero alternating
sum on the permutation orbit. To see this, since $H\ge3$, either $T$
or its complement contains two distinct positions. Swapping those
positions preserves the monomial's value and reverses the permutation
sign, pairing all summands with their negatives. Linearity therefore
gives
\[
 \sum_{\pi\in\mathfrak S_H}\operatorname{sgn}(\pi)
 P(\lambda_{\pi(1)},\ldots,\lambda_{\pi(H)})=0.
\]
The proposed strict signs would make every summand positive, a
contradiction. Thus even fully asymmetric queries cannot read an
arbitrary labeling of the factorial code. This particular obstruction
does not bound its shattering dimension, since a shattering construction
need not use permutation parity as one of its readable bits.

\paragraph{Arbitrary assignment costs do have superlinear shattering.}
The following finite construction isolates what the reciprocal transfer
would have to supply. Let $L=2^k$ be sufficiently large, set $H=L^2$
and $m=H/2$. There exist $m$ colorings
$c_g:\{1,\ldots,H\}\to\{1,\ldots,L\}$ such that every prescription
of one color for each $g$ can be realized by distinct values $j_g$.
Consequently, arbitrary assignment-cost queries on permutations can
shatter
\[
 mk=\frac14H\log_2H
\]
bits. This claim does not restrict the costs to reciprocal functions.

Here is a full existence proof. Independently for each $g$, partition
the $H$ values uniformly into $L$ labeled bins of size $L$.
A selection of one bin per slot fails to have distinct representatives
only if some $t$ selected bins have union contained in a set $U$ of
size $t-1$, by Hall's theorem. Necessarily $t\ge L+1$, since every
bin has size $L$. For fixed $t$ slots, selected bin labels, and $U$,
the probability that all bins lie in $U$ is
\[
 \left(\frac{\binom{t-1}{L}}{\binom HL}\right)^t
 \le\left(\frac tH\right)^{Lt}.
\]
Independence holds across slots; within a slot only one prescribed bin
is being tested. Taking a union bound over slots, labels, and $U$
gives failure probability at most
\[
 \sum_{t=L+1}^{m}\binom mt\binom H{t-1}L^t
                  \left(\frac tH\right)^{Lt}.
\]
Using $\binom mt\le(eH/t)^t$ and
$\binom H{t-1}\le(2eH/t)^t$, each summand is at most
\[
 \left[2e^2L\left(\frac tH\right)^{L-2}\right]^t
 \le\left[8e^2L\,2^{-L}\right]^t,
\]
because $t\le H/2$. For all sufficiently large $L$ the bracket is
at most $1/2$, and the sum is less than one. A choice of partitions
with no Hall failure therefore exists.

For each slot $g\le m$ and bit index $\ell\le k$, define a query
on permutations by the corresponding bit of $c_g(\pi(g))$, represented
as a sign in $\{-1,1\}$. It is an assignment-cost sum
$\sum_h C_{g,\ell}(h,\pi(h))$ with all cost rows zero except $h=g$.
Given any $mk$ desired bits, choose their $m$ desired colors. The
proved distinct-representative property assigns distinct values to
the first $m$ slots; extend this injection to a permutation of all
$H$ values. Every desired strict sign is then realized.

Even monotonicity of the cost rows alone does not block this example.
For any fixed query cost matrix $C$ and ordered distinct real values
$\lambda_1<\cdots<\lambda_H$, choose $M$ sufficiently large that
$C(h,j)+M\lambda_j$ is strictly increasing in $j$ for every $h$.
Adding this same column function to every row changes an assignment
sum by the permutation-independent constant $M\sum_j\lambda_j$,
which a bias removes exactly. Thus all costs can be made increasing
without changing any query sign. However, these arbitrary increasing
cost vectors need not have the one-pole affine-ratio form. The
one-slot bound above proves that their transfer to that form actually
fails for this construction. No reciprocal shattering conclusion
follows from the assignment-cost theorem.

\paragraph{The near-endpoint pole proposal and its exact budget.}
For isolated scalar values $\lambda_j=\varepsilon-R^{-j}$, the kernel
\[
 \frac{t}{\varepsilon+t-\lambda_j}
 =\frac{t}{t+R^{-j}}
\]
indeed gives a rank threshold after comparison with $1/2$.
However, its direct use in the reverse-orientation chart has a
whole-cube constraint. Write
\[
 B(u,v)=C-\sum_{i=1}^r a_i u_i-\sum_{j=1}^S b_j v_j,
 \qquad a_i,b_j>0,\qquad C>\sum_i a_i+\sum_j b_j.
\]
The row baseline $d_i=C-a_i$ necessarily satisfies
\[
 d_i>\sum_{k\ne i}a_k+\sum_jb_j>\sum_jb_j
\]
when $r>1$, and $d_i>\sum_jb_j$ also when $r=1$.
In particular, if every independently coded column has
$b_j>\varepsilon-\eta$, then $d_i>S(\varepsilon-\eta)$.
For two or more such columns this excludes placing all row poles
arbitrarily close to $\varepsilon$.

Even with a single column there is a row budget. Put
$C=\varepsilon+T$, $d_i=\varepsilon+t_i$, with $T,t_i>0$,
and let its parameter be $b$.
Then $a_i=T-t_i>0$, and whole-cube positivity gives exactly
\[
 \sum_i(T-t_i)<T+\varepsilon-b.
\]
Thus, for $0<\theta<1$, the number of rows with $t_i\le\theta T$
is strictly less than
\[
 \frac{1+(\varepsilon-b)/T}{1-\theta}.
\]
Many freely placed geometric thresholds much smaller than the largest
threshold therefore do not fit in one head when $b$ approaches
$\varepsilon$. These inequalities refute only that direct placement
of near-endpoint kernels. They do not rule out many-head cancellation,
other parameter encodings, or superlinear reciprocal shattering.

The open step for this local route is to construct sufficiently many
asymmetric reciprocal queries whose strict sign vectors on, or near,
these sheets cover all binary words of length $cH\log H$ for some
constant $c>0$. No such statement is assumed or proved here. A generic
upper bound on the number of sign patterns would not supply it.

\end{document}
